\documentclass[10pt,a4paper]{amsart}
\usepackage[margin=19mm]{geometry}
\usepackage{amsmath,amssymb,mathtools}
\usepackage[scr=rsfs,cal=euler]{mathalfa}
\usepackage{xfrac}
\usepackage[unicode,hidelinks,bookmarks=false]{hyperref}
\newcommand{\ie}{i.e.\ }
\newcommand{\cf}{cf.\ }
\newcommand{\resp}{respectively\ }
\newcommand{\Subsection}{\subsection}
\providecommand{\nobreakdash}{\nobreak\hbox{-}}
\setlength{\emergencystretch}{2em}
\multlinegap0pt
\usepackage{amssymb}
\usepackage{mathtools}
\usepackage{csquotes}
\usepackage{algorithm}\usepackage{algpseudocode}
\usepackage{tikz}
\usepackage{enumerate}
\newtheorem{lem}{Lemma}
\title{On the number of rational power factors in a finite word}
\author{Shuo Li, Yuan Song}
\date{Source: arXiv:2605.14955v1 (CC BY 4.0)}
\begin{document}
\maketitle

\noindent\emph{Source and licence.} On the number of rational power factors in a finite word. Version arXiv:2605.14955v1 (CC BY 4.0), distributed by arXiv under the Creative Commons Attribution 4.0 International licence, http://creativecommons.org/licenses/by/4.0/, as recorded in the arXiv licence metadata for this exact version and shown on the abstract page https://arxiv.org/abs/2605.14955v1. Full licence notice: \texttt{licenses/arxiv-2605.14955-CC-BY-4.0.txt}.\par\smallskip

\noindent\textit{Editorial scope.} Counted unit: the article's lem lem:packing-function (source0.tex lines 599--620). The article's complete original proof of it is delivered with it (source0.tex lines 622--738, one \texttt{proof} environment of 117 lines). No mathematics is written, repaired or re-derived: the statement and the proof are the article's own lines, in the article's own notation, and no retained line is substituted or re-flowed.  No further source-local context is needed: every cross-reference of every retained line resolves inside the two delivered spans.  Nothing was omitted from any retained span, so the package carries no omission record.  No retained line contains a citation, so the wrapper carries no bibliography and no bibliography entry.  No retained line contains a figure environment, an image-inclusion command or drawing code, so no image is delivered and the package declares no asset dependency.  Editorial formatting only, in addition: an AMS \texttt{amsart} preamble that restates the article's own declarations and macros used by the retained lines, namely the environments \texttt{lem} on separate counters, together with the standard packages those lines need.  The retained environments print in this document as Lemma 1 in order of appearance, whereas the article prints the same items with the numbering of its own full document.  The complete native source of the article as distributed in the arXiv e-print for this exact version is bundled unchanged as \texttt{sources/200626/source0.tex}.
\begin{lem}\label{lem:packing-function}
Let \(0<s\le 1/2\) and \(s\le B\le 2s\). Then
\[
\mathcal P(s,B)\le s^2V(B/s),
\]
where, for \(1\le T\le 2\),
\[
V(T):=
\begin{cases}
\dfrac{(2-T)(3T-2)}{4},
& 1\le T\le \dfrac43,\\[8pt]
\dfrac32T^2-4T+3,
& \dfrac43\le T\le \dfrac32,\\[8pt]
\dfrac12(3-T)(T-1),
& \dfrac32\le T\le 2.
\end{cases}
\]
In particular,
\[
\mathcal P\left(\frac12,1\right)\le \frac18.
\]
\end{lem}

\begin{proof}
Let $T:=\frac{B}{s}$, then \(1\le T\le 2\).

Fix one pair \((x,y)\) with $x+y\le s$. Write
\[
y=rs,
\qquad 0\le r\le 1.
\]
For a fixed \(y\), the largest possible value of \(F(x,y)\), under the condition
\(0\le x\le s-y\), is $s^2\varphi(r)$, where
\[
\varphi(r):=
\begin{cases}
\dfrac12r^2, & 0\le r\le \dfrac12,\\[6pt]
\dfrac12(1-r)(3r-1), & \dfrac12\le r\le 1.
\end{cases}
\]
Indeed, if \(r\le 1/2\), then one can choose \(x\ge y\), and the maximum is
\[
\frac12y^2=s^2\frac12r^2.
\]
If \(r\ge 1/2\), then \(x\le s-y\le y\), and
\[
F(x,y)=xy-\frac12x^2
\]
is increasing in \(x\) on \(0\le x\le s-y\).  Hence the maximum is attained at
\(x=s-y\), and equals
\[
(s-y)y-\frac12(s-y)^2
=
s^2\frac12(1-r)(3r-1).
\]
Therefore, for every admissible family \((x_j,y_j)\), if $r_j:=\frac{y_j}{s}$, then
\[
F(x_j,y_j)\le s^2\varphi(r_j),\quad\sum_j r_j\le T.
\]
It remains to maximize $\sum_j\varphi(r_j)$ under the constraints
\[
0\le r_j\le 1,
\qquad
\sum_j r_j\le T.
\]
Let $r=T-\sum_jr_j$, then $\max \sum_j\varphi(r_j) \leq \max (\sum_j\varphi(r_j)+\varphi(r))$. We can now compute $ \max \sum_j\varphi(r_j)$ under the restrictions:
\[
0\le r_j\le 1,
\qquad
\sum_j r_j= T.
\]
We remark that for any pair $0 \le r_i, r_j<\frac{1}{2}$ with $r_i \ne r_j$, if $r_i+r_j \leq \frac{1}{2}$, then $$\varphi(r_i)+\varphi(r_j) \leq \varphi(r_i+r_j);$$ if $r_i+r_j > \frac{1}{2}$, then $$\varphi(r_i)+\varphi(r_j) \leq \varphi(\frac{1}{2})+ \varphi(r_i+r_j-\frac{1}{2}).$$ Thus, one can reduce the number of variables in objective function $\sum_j\varphi(r_j)$ into at most $4$. That is, $\max \sum_j\varphi(r_j) \le \max \sum_{j=1}^k\varphi(r'_j)$ under the restrictions: 
\[
0\le r'_1 \le r'_2 \le \ldots \le r'_k\le 1
,\qquad
\sum_{j=1}^k r'_j= T
,\qquad r'_2 \ge \frac{1}{2}.
\]
Moreover, since $\varphi$ is concave in $[\frac{1}{2},1]$, 
$$\sum_{j=2}^k\varphi(r'_j) \leq (k-1)\varphi(\frac{1}{k-1}\sum_{j=2}^kr'_j).$$
Hence, $$\sum_{j=1}^k\varphi(r'_j) \leq \varphi(r'_1)+(k-1)\varphi(\frac{T-r'_1}{k-1}).$$
\iffalse
The function \(\varphi\) is convex on \([0,1/2]\) and concave on \([1/2,1]\).
Thus, in an extremal family, there is at most one entry in \((0,1/2)\).  The
remaining nonzero entries lie in \([1/2,1]\).  Since \(T\le 2\), there are at
most three entries strictly larger than \(1/2\).  The endpoint configuration
with four entries in \([1/2,1]\) can occur only when \(T=2\), in which case all
four entries are equal to \(1/2\), and the value is
\[
4\varphi(1/2)=\frac12=V(2).
\]

It remains to check the cases with \(k=1,2,3\) entries in \([1/2,1]\), together
with possibly one entry \(z\in[0,1/2]\).  For fixed \(k\) and \(z\), concavity on
\([1/2,1]\) implies that the \(k\) large entries are equal.  Hence we only need
to maximize
\[
k\varphi\left(\frac{T-z}{k}\right)+\frac12z^2,
\]
where
\[
k=1,2,3,\qquad
0\le z\le\frac12,
\qquad
\frac12\le\frac{T-z}{k}\le 1.
\]
\fi
For $k=2,3,4$, a direct one-variable calculation gives the following three possible maxima:
\[
2\varphi(T/2)
\quad
\text{for }
1\le T\le \frac43,
\]
\[
2\varphi(2-T)+\frac12(3T-4)^2
\quad
\text{for }
\frac43\le T\le \frac32,
\]
and
\[
3\varphi(T/3)
\quad
\text{for }
\frac32\le T\le 2.
\]
These are exactly
\[
V(T)=
\begin{cases}
\dfrac{(2-T)(3T-2)}{4},
& 1\le T\le \dfrac43,\\[8pt]
\dfrac32T^2-4T+3,
& \dfrac43\le T\le \dfrac32,\\[8pt]
\dfrac12(3-T)(T-1),
& \dfrac32\le T\le 2.
\end{cases}
\]\qed
\end{proof}

\end{document}
